\PassOptionsToPackage{table}{xcolor}
\documentclass[10pt,aspectratio=169]{beamer}
\newcommand{\LectureNumber}{02}
\input{course-theme.tex}
\title{Source files and execution}
\subtitle{CSC103 Programming Fundamentals / Lecture 02}
\author{Dr. Muhammad Farhan}
\institute{Associate Professor of Computer Science\\Department of Computer Science\\COMSATS University Islamabad\\Sahiwal Campus}
\date{Fall 2026}
\begin{document}
\begin{frame}\titlepage\end{frame}

\begin{frame}[fragile,t]{The problem for this lecture}
\begin{columns}[T,onlytextwidth]\begin{column}{0.60\textwidth}
Create Welcome.java. Print your name and the result of 9 * 4 on separate lines. Compile and execute it.
\end{column}\begin{column}{0.35\textwidth}\small
Create a source file whose public class matches its filename.
\end{column}\end{columns}
\note{Introduce the target problem before explaining the tools. Revisit it in the guided solution later.\par Syllabus reading: Deitel Ch 1. CLO-1.}
\end{frame}

\begin{frame}[fragile,t]{Learning objectives}
\begin{columns}[T,onlytextwidth]\begin{column}{0.60\textwidth}
\begin{itemize}
\item Create and compile a Java source file
\item Explain the Java execution pipeline
\item Classify syntax, runtime and logic errors
\end{itemize}
\end{column}\begin{column}{0.35\textwidth}\small
CLO-1

A first Java source file\par Compilation and execution\par Three error categories\par Files and working directories
\end{column}\end{columns}
\note{Explain how each objective supports the opening problem. The final questions check these objectives.\par Syllabus reading: Deitel Ch 1. CLO-1.}
\end{frame}

\begin{frame}[fragile,t]{A first Java source file}
\begin{itemize}
\item A public class name matches its .java filename.
\item The conventional entry point is public static void main(String[] args).
\end{itemize}
\vspace{1mm}\begin{block}{Example}
\begin{lstlisting}
public class Hello {
  public static void main(String[] args) {
    System.out.println("Hello, Java");
  }
}
\end{lstlisting}
\end{block}
\note{Save as Hello.java. Explain that braces group the class and method bodies. Formal treatment of methods comes later.\par Syllabus reading: Deitel Ch 1. CLO-1.}
\end{frame}

\begin{frame}[fragile,t]{Compilation and execution}
\begin{itemize}
\item javac translates source code into JVM bytecode.
\item java starts the JVM, which loads and links classes, then executes the program.
\end{itemize}
\vspace{1mm}\begin{block}{Example}
\begin{lstlisting}
javac Hello.java
java Hello

Hello.java   source
Hello.class  bytecode
\end{lstlisting}
\end{block}
\note{In Java, linking involves loading, verification, preparation and resolution. There is normally no separate user linker command as in a C toolchain.\par Syllabus reading: Deitel Ch 1. CLO-1.}
\end{frame}

\begin{frame}[fragile,t]{Three error categories}
\begin{itemize}
\item Syntax and type errors prevent successful compilation.
\item Runtime exceptions interrupt execution. Logic errors produce an incorrect result even when execution finishes.
\end{itemize}
\vspace{1mm}\begin{block}{Example}
\begin{lstlisting}
Syntax: missing ;
Runtime: Integer.parseInt("cat")
Logic: area = length + width
\end{lstlisting}
\end{block}
\note{Use actual compiler diagnostics in a demo. Read the first meaningful diagnostic first because one missing symbol can cause several later messages.\par Syllabus reading: Deitel Ch 1. CLO-1.}
\end{frame}

\begin{frame}[fragile,t]{Files and working directories}
\begin{itemize}
\item A terminal command runs relative to its working directory.
\item An editor or IDE helps save, build and run code, but the Java compiler still enforces language rules.
\end{itemize}
\vspace{1mm}\begin{block}{Example}
\begin{lstlisting}
1. Create a course folder.
2. Save Hello.java as plain text.
3. Open a terminal in that folder.
4. Compile, then run.
\end{lstlisting}
\end{block}
\note{Check for hidden .txt extensions. On systems with multiple JDKs, compare java -version and javac -version. These examples target Java 17 language features.\par Syllabus reading: Deitel Ch 1. CLO-1.}
\end{frame}

\begin{frame}[fragile,t]{What compilation checks}
\begin{itemize}
\item The compiler checks syntax and type consistency before producing bytecode.
\item It cannot generally decide whether your formula answers the problem correctly.
\item A successful build therefore needs a separate execution and result check.
\end{itemize}
\vspace{1mm}\begin{block}{Example}
\begin{lstlisting}
int n = "five";  // incompatible types
int area = 5 + 3; // legal, but wrong for a 5 by 3 area
\end{lstlisting}
\end{block}
\note{Explain the mechanism using the displayed example. The compiler checks syntax and type consistency before producing bytecode. It cannot generally decide whether your formula answers the problem correctly. A successful build therefore needs a separate execution and result check.\par Syllabus reading: Deitel Ch 1. CLO-1.}
\end{frame}

\begin{frame}[fragile,t]{Reading an error message}
\begin{itemize}
\item A diagnostic normally identifies a file, line and kind of problem.
\item Start with the earliest relevant error because one mistake can trigger several messages.
\item Recompile after a small correction to see which messages remain.
\end{itemize}
\vspace{1mm}\begin{block}{Example}
\begin{lstlisting}
Missing semicolon near a println statement:
Check that statement and the previous line.
\end{lstlisting}
\end{block}
\note{Explain the mechanism using the displayed example. A diagnostic normally identifies a file, line and kind of problem. Start with the earliest relevant error because one mistake can trigger several messages. Recompile after a small correction to see which messages remain.\par Syllabus reading: Deitel Ch 1. CLO-1.}
\end{frame}


\begin{frame}[fragile,t]{The Java execution pipeline}
\begin{center}\begin{tikzpicture}
\node[box] (n0) at (0.0,0) {Source .java};
\node[box] (n1) at (3.45,0) {javac compiler};
\node[box] (n2) at (6.9,0) {Bytecode .class};
\node[alt] (n3) at (10.350000000000001,0) {JVM execution};
\draw[arr] (n0) -- (n1);
\draw[arr] (n1) -- (n2);
\draw[arr] (n2) -- (n3);
\end{tikzpicture}\end{center}
\begin{block}{What to notice}Compilation and execution are different stages; each needs its own check.\end{block}
\note{Trace each labelled connection. Compilation and execution are different stages; each needs its own check.}
\end{frame}
\begin{frame}[fragile,t]{Key distinctions}
\small\renewcommand{\arraystretch}{1.5}
\rowcolors{2}{pale}{white}
\begin{tabularx}{\textwidth}{>{\raggedright\arraybackslash}X>{\raggedright\arraybackslash}X>{\raggedright\arraybackslash}X}
\rowcolor{navy}
\color{white}\bfseries Stage & \color{white}\bfseries Input & \color{white}\bfseries Result\\
Edit & Plain-text source & Welcome.java\\
Compile & javac Welcome.java & Welcome.class or diagnostics\\
Load and link & java Welcome & JVM prepares classes\\
Execute & main statements & Console output\\
\end{tabularx}
\vspace{3mm}\footnotesize These distinctions determine which operation or control structure fits the problem.
\note{\par Syllabus reading: Deitel Ch 1. CLO-1.}
\end{frame}

\begin{frame}[fragile,t]{First example: execution}
\begin{lstlisting}
System.out.println("Welcome to CSC103");
System.out.println(7 + 5);
\end{lstlisting}
\note{This is the main body from Java\_Examples/Lecture02.java. The source file includes the complete class. Predict the result before the next slide.\par Syllabus reading: Deitel Ch 1. CLO-1.}
\end{frame}

\begin{frame}[fragile,t]{First example: result}
\begin{columns}[T,onlytextwidth]\begin{column}{0.60\textwidth}
Welcome to CSC103\par 12\par The first statement prints text. The second prints an arithmetic result.
\end{column}\begin{column}{0.35\textwidth}\small
Three error categories

Files and working directories
\end{column}\end{columns}
\note{Expected output and interpretation: Welcome to CSC103
12
The first statement prints text. The second prints an arithmetic result.\par Syllabus reading: Deitel Ch 1. CLO-1.}
\end{frame}

\begin{frame}[fragile,t]{Guided practice}
\begin{columns}[T,onlytextwidth]\begin{column}{0.60\textwidth}
Create Welcome.java. Print your name and the result of 9 * 4 on separate lines. Compile and execute it.
\end{column}\begin{column}{0.35\textwidth}\small
Attempt the solution before continuing.

Use the definitions and examples from this lecture.
\end{column}\end{columns}
\note{Give students 8 minutes to attempt the problem. The following slides provide a complete solution. Original solution guidance: The public class must be Welcome. Place both println statements inside main. Expected second line: 36.\par Syllabus reading: Deitel Ch 1. CLO-1.}
\end{frame}

\begin{frame}[fragile,t]{Solution strategy}
\begin{columns}[T,onlytextwidth]\begin{column}{0.60\textwidth}
\begin{itemize}
\item Create a source file whose public class matches its filename.
\item Put the two output statements inside main.
\item Compile the file, then run the class without a .class extension.
\end{itemize}
\end{column}\begin{column}{0.35\textwidth}\small
The next program uses fixed inputs so its result can be reproduced.
\end{column}\end{columns}
\note{Discuss why each step is needed. State the input assumptions before executing the program.\par Syllabus reading: Deitel Ch 1. CLO-1.}
\end{frame}

\begin{frame}[fragile,t]{Complete solution}
\begin{lstlisting}
public class Solution02 {
  public static void main(String[] args) throws Exception {
    System.out.println("Amina");
    System.out.println(9 * 4);
  }
}
\end{lstlisting}
\note{Complete runnable file: Java\_Examples/Solution02.java. Inputs are fixed in the program.  \par Syllabus reading: Deitel Ch 1. CLO-1.}
\end{frame}

\begin{frame}[fragile,t]{Execution trace}
\small\renewcommand{\arraystretch}{1.5}
\rowcolors{2}{pale}{white}
\begin{tabularx}{\textwidth}{>{\raggedright\arraybackslash}X>{\raggedright\arraybackslash}X>{\raggedright\arraybackslash}X}
\rowcolor{navy}
\color{white}\bfseries Statement & \color{white}\bfseries Evaluation & \color{white}\bfseries Console line\\
println("Amina") & String literal & Amina\\
println(9 * 4) & 9 multiplied by 4 & 36\\
End of main & Normal return & No additional line\\
\end{tabularx}
\vspace{3mm}\footnotesize Follow the changing state in execution order.
\note{Manually derived trace for the complete solution. Create a source file whose public class matches its filename. Put the two output statements inside main. Compile the file, then run the class without a .class extension.\par Syllabus reading: Deitel Ch 1. CLO-1.}
\end{frame}

\begin{frame}[fragile,t]{Solution result}
\begin{columns}[T,onlytextwidth]\begin{column}{0.60\textwidth}
Amina\par 36
\end{column}\begin{column}{0.35\textwidth}\small
Why this result follows

Compile the file, then run the class without a .class extension.
\end{column}\end{columns}
\note{Expected result: Amina
36\par Syllabus reading: Deitel Ch 1. CLO-1.}
\end{frame}

\begin{frame}[fragile,t]{A common incorrect approach}
File name: Start.java\par public class Begin \{ \}
\vspace{1mm}\begin{block}{Example}
\begin{lstlisting}
Rename the file Begin.java or the public class Start. Java requires the public top-level class and filename to match.
\end{lstlisting}
\end{block}
\note{Ask students to predict the failure before discussing the explanation. The right side gives the correction.\par Syllabus reading: Deitel Ch 1. CLO-1.}
\end{frame}

\begin{frame}[fragile,t]{Boundary and failure checks}
\begin{columns}[T,onlytextwidth]\begin{column}{0.60\textwidth}
Check the stated input assumptions.

Create one example of each error category. Record the error message or wrong result and its correction.
\end{column}\begin{column}{0.35\textwidth}\small
The public class must be Welcome. Place both println statements inside main. Expected second line: 36.
\end{column}\end{columns}
\note{Use the specification to derive each expected result. Exercise solution: The public class must be Welcome. Place both println statements inside main. Expected second line: 36.\par Syllabus reading: Deitel Ch 1. CLO-1.}
\end{frame}

\begin{frame}[fragile,t]{Check your understanding}
\begin{columns}[T,onlytextwidth]\begin{column}{0.60\textwidth}
Does java Hello compile Hello.java in the demonstrated two-step workflow? What file does javac create?
\end{column}\begin{column}{0.35\textwidth}\small
Write a prediction and explain the rule that supports it.
\end{column}\end{columns}
\note{Answer: No. java Hello runs the class. javac Hello.java creates Hello.class. Java also has a separate source-file launch mode, outside this demonstration.\par Syllabus reading: Deitel Ch 1. CLO-1.}
\end{frame}

\begin{frame}[fragile,t]{Answer and explanation}
\begin{columns}[T,onlytextwidth]\begin{column}{0.60\textwidth}
No. java Hello runs the class. javac Hello.java creates Hello.class. Java also has a separate source-file launch mode, outside this demonstration.
\end{column}\begin{column}{0.35\textwidth}\small
Rename the file Begin.java or the public class Start. Java requires the public top-level class and filename to match.
\end{column}\end{columns}
\note{Reconnect the answer to the relevant definition rather than treating it as a fact to memorise.\par Syllabus reading: Deitel Ch 1. CLO-1.}
\end{frame}


\begin{frame}[fragile,t]{Compare cases and predict outcomes}
\small\renewcommand{\arraystretch}{1.5}\rowcolors{2}{pale}{white}
\begin{tabularx}{\textwidth}{>{\raggedright\arraybackslash}X>{\raggedright\arraybackslash}X>{\raggedright\arraybackslash}X}\rowcolor{navy}
\color{white}\bfseries Observation & \color{white}\bfseries Stage & \color{white}\bfseries Next check\\
Missing semicolon & Compilation & Inspect this and the preceding line\\
Class not found & Launch / loading & Class name and classpath\\
Wrong printed total & Execution / logic & Formula and test input\\
\end{tabularx}\par\vspace{3mm}\begin{exampleblock}{Discuss}Explain each expected result before running the example. Which case would reveal a wrong assumption?\end{exampleblock}
\note{Read the first two columns and invite predictions before discussing the outcome.}
\end{frame}

\begin{frame}[fragile,t]{Apply the idea}
\begin{alertblock}{Independent practice}A public class named Receipt is saved as Bill.java. After renaming it, the program prints 5 + 3 when the required area is 5 by 3. Identify both corrections and the commands.\end{alertblock}\vspace{3mm}\begin{enumerate}\item Work through the input by hand.\item Write or correct the program.\item Justify the expected result.\end{enumerate}\note{Allow 3–5 minutes, then compare answers in pairs. A public class named Receipt is saved as Bill.java. After renaming it, the program prints 5 + 3 when the required area is 5 by 3. Identify both corrections and the commands.}
\end{frame}

\begin{frame}[fragile,t]{Solution and reasoning}
\begin{lstlisting}
// Save as Receipt.java
public class Receipt {
  public static void main(String[] args) {
    System.out.println(5 * 3);
  }
}
// javac Receipt.java
// java Receipt
\end{lstlisting}\vspace{1mm}\small\begin{itemize}
\item The public class and source filename must agree.
\item Changing + to * fixes the computation; successful compilation alone cannot establish that.
\item Expected output: 15.
\end{itemize}\note{Answer: The public class and source filename must agree. Changing + to * fixes the computation; successful compilation alone cannot establish that. Expected output: 15.}
\end{frame}

\begin{frame}[fragile,t]{Java program phases and runtime roles}
\begin{itemize}
\item javac checks source and emits class{-}file bytecode; the java launcher starts execution through a JVM.
\item Loading, verification and execution are distinct runtime responsibilities. JIT compilation works on bytecode, not Java source.
\item Use a JDK for development. Modern JDK layouts no longer use the old rt.jar structure shown in the source.
\end{itemize}
\vspace{3mm}\footnotesize\textcolor{accent}{Source: supplied ISB campus slides. See speaker notes and ISB\_Source\_Map.md.}
\note{Adapted from the supplied ISB campus folder: Lecture{-}1.KI.pptx, slides 18, 25, 26, 27, 28, 29. Corrected the legacy rt.jar/JRE description and source{-}versus{-}bytecode wording. Build commands use the supplied IsbLecture02 filename.}
\end{frame}

\begin{frame}[fragile,t]{Worked problem: Java program phases and runtime roles}
\begin{block}{Task}Build and run a small program, then identify which phase detects a missing semicolon and which phase prints the greeting.\end{block}
\small Input: Values are fixed in the program for a reproducible demonstration.\par\vspace{3mm}\begin{exampleblock}{Before running}Predict the output and identify the variables that change.\end{exampleblock}
\note{Adapted from the supplied ISB campus folder: Lecture{-}1.KI.pptx, slides 18, 25, 26, 27, 28, 29. Corrected the legacy rt.jar/JRE description and source{-}versus{-}bytecode wording. Build commands use the supplied IsbLecture02 filename.}
\end{frame}

\begin{frame}[fragile,t]{Complete ISB example (1/1)}
\begin{lstlisting}
public class IsbLecture02 {
  public static void main(String[] args) throws Exception {
    System.out.println("Hello, World");
  }

}
\end{lstlisting}
\note{Adapted from the supplied ISB campus folder: Lecture{-}1.KI.pptx, slides 18, 25, 26, 27, 28, 29. Corrected the legacy rt.jar/JRE description and source{-}versus{-}bytecode wording. Build commands use the supplied IsbLecture02 filename. Complete file: ISB\_Java\_Examples/IsbLecture02.java.}
\end{frame}

\begin{frame}[fragile,t]{Worked trace and expected result}
\small\renewcommand{\arraystretch}{1.4}\rowcolors{2}{pale}{white}
\begin{tabularx}{\textwidth}{>{\raggedright\arraybackslash}X>{\raggedright\arraybackslash}X>{\raggedright\arraybackslash}X}\rowcolor{navy}
\color{white}\bfseries Stage & \color{white}\bfseries Artifact or action & \color{white}\bfseries Evidence\\
Compile & javac IsbLecture02.java & Class file or diagnostic\\
Load and verify & Prepare class bytecode & Valid class structure required\\
Execute & Run main & Hello, World\\
\end{tabularx}\par\vspace{2mm}
\begin{block}{Expected console output}\ttfamily\footnotesize Hello, World\end{block}
\note{Adapted from the supplied ISB campus folder: Lecture{-}1.KI.pptx, slides 18, 25, 26, 27, 28, 29. Corrected the legacy rt.jar/JRE description and source{-}versus{-}bytecode wording. Build commands use the supplied IsbLecture02 filename. Trace and expected values independently worked for this edition.}
\end{frame}

\begin{frame}[fragile,t]{Extend the ISB example}
\begin{alertblock}{Practice}Does bytecode verification prove that the formula in a program is correct?\end{alertblock}\vspace{3mm}\begin{exampleblock}{Explain your reasoning}No. Verification checks structural and type{-}related constraints; a valid program can still compute the wrong result.\end{exampleblock}
\note{Adapted from the supplied ISB campus folder: Lecture{-}1.KI.pptx, slides 18, 25, 26, 27, 28, 29. Corrected the legacy rt.jar/JRE description and source{-}versus{-}bytecode wording. Build commands use the supplied IsbLecture02 filename. Answer: No. Verification checks structural and type{-}related constraints; a valid program can still compute the wrong result.}
\end{frame}
\begin{frame}[fragile,t]{Lecture summary}
\begin{columns}[T,onlytextwidth]\begin{column}{0.60\textwidth}
\begin{itemize}
\item and compile a Java source file
\item the Java execution pipeline
\item syntax, runtime and logic errors
\end{itemize}
\end{column}\begin{column}{0.35\textwidth}\small
\begin{itemize}
\item Create a source file whose public class matches its filename.
\item Put the two output statements inside main.
\item Compile the file, then run the class without a .class extension.
\end{itemize}
\end{column}\end{columns}
\note{Ask students to give one concrete example for the concepts listed. Use the worked solution to connect the topics.\par Syllabus reading: Deitel Ch 1. CLO-1.}
\end{frame}

\begin{frame}[fragile,t]{After-class practice}
\begin{columns}[T,onlytextwidth]\begin{column}{0.60\textwidth}
Create one example of each error category. Record the error message or wrong result and its correction.
\end{column}\begin{column}{0.35\textwidth}\small
Syllabus reading\par Deitel Ch 1

Next lecture: Java syntax and program style
\end{column}\end{columns}
\note{Reading labels follow the supplied outline. Check topic headings against the available textbook edition. Require a program and expected results for the practice task.\par Syllabus reading: Deitel Ch 1. CLO-1.}
\end{frame}
\end{document}
